\chapter{Architecture}
\label{ch:architettura_en}

\section{Stack}

\begin{itemize}[leftmargin=*]
  \item \textbf{Backend}: Python 3.11+, FastAPI, SQLAlchemy 2.0,
        Pydantic 2, uvicorn, OR-Tools, NumPy, scikit-learn, Faker.
  \item \textbf{Frontend}: SvelteKit (Svelte 4), Vite 5, Tailwind
        CSS, plain JavaScript.
  \item \textbf{DB}: SQLite via SQLAlchemy. Idempotent migrations
        in code (no Alembic).
  \item \textbf{Engine I/O}: Python pickle;
        \texttt{engine\_io.py} converts between DB and dict for
        the solver.
\end{itemize}

\section{Three-tier overview}

The system has a CP-SAT solver core (\texttt{engine/}), a
FastAPI + SQLite backend (\texttt{webui/backend/}) and a
SvelteKit frontend (\texttt{webui/frontend/}). The solver runs
as a library imported by the backend; pipelines are exposed as
async runs via \texttt{/api/optimize/*}. The frontend talks
to the backend over JSON; long-running solver runs are tracked
via the \texttt{runs} table polled at 1\,Hz from the UI.

Solver pipelines (in \texttt{engine/}):
\begin{itemize}
\item \texttt{cpsat\_v2\_assignment.py}: Phase A (teacher to
  class assignment).
\item \texttt{cpsat\_v2\_timetable.py}: Phase A (timetabling,
  \texttt{day\_count}) + Phase B (per-day slot placement).
  Native locks + C1/C2/C3 constraints.
\item \texttt{decomposition\_temporal.py}: 6-day parallel
  decomposition.
\item \texttt{decomposition\_spectral\_v2.py}, \texttt{curriculum}
  / \texttt{metis} variants: cluster classes, solve
  sub-problems, ricucitura.
\item \texttt{column\_generation.py}: master LP + diversified
  pattern enrichment + completion fallback. Mode
  \texttt{branch-and-price} fully wired with all scalability
  techniques (RF tree, dual stabilization, column management,
  parallel pricing).
\item \texttt{metaheuristics.py}, \texttt{alns.py},
  \texttt{vns.py}, \texttt{lagrangian.py}: post-processing on a
  HARD-feasible seed solution.
\end{itemize}

\section{Solver architecture: from DSL to CP-SAT}
\label{sec:arch-dsl-cpsat-en}

Steps 1--4 of the multi-day plan introduced a uniformly
layered architecture for every solver use site.

\begin{verbatim}
+-----------------------------------------------------------+
| UI (SvelteKit)                                            |
|   - Constraints tab: matrices + DSL + plessi              |
|   - Workflow tab: pipeline (mode, granularity)            |
+-----------------+----------------------------+------------+
                  |                              |
        POST /api/constraints/general     POST /api/optimize/*
                  |                              |
                  v                              v
+-----------------+----------------+   +---------+---------+
| dsl_parser (general_dsl.py)      |   | router            |
|   parse() -> AST                  |   | run_async         |
+-----------------+----------------+   +---------+---------+
                  |                              |
                  v                              v
+-----------------+----------------+   +---------+---------+
| dsl_translator                   |   | engine pipeline   |
|   special-purpose tables         |   | (mono / decomp /  |
|   -> canonical DSL strings       |   |  CG / BP)         |
+-----------------+----------------+   +---------+---------+
                  |                              |
                  +------------- merged ---------+
                                  |
                                  v
+--------------------------------+--------------------------+
| dsl_to_cpsat.DSLConstraintCompiler                         |
|   AST + slot dict + plessi_data                            |
|   -> add_assertion() / add_implication() on CpModel        |
+--------------------------------+--------------------------+
                                 |
              +------------------+------------------+
              v                  v                  v
       MonolithicSolver   PhaseBDaySolver   BP Pricer (9 grans)
       (cpsat_v2_*)       (legacy +         + RF nodes
       seed via DSL       DSL augment)
              |                  |                  |
              +------------------+------------------+
                                 |
                                 v
                         solution (Lessons)
                                 |
                                 v
                          DSL post-hoc evaluator
                          (HARD: feasibility check)
                          (SOFT: contributes to
                                  global score)
\end{verbatim}

Key properties:
\begin{itemize}
\item \textbf{One parser, one AST}. Every constraint source
  (matrices, DSL UI, plesso rules, seeded legacy HARDs) is
  unified into canonical DSL strings parsed by
  \texttt{general\_dsl.py}.
\item \textbf{One compiler}. \texttt{DSLConstraintCompiler} is
  the single gate to CP-SAT. Mono solver, BP pricers,
  PhaseBDaySolver, Ryan-Foster nodes -- they all apply the same
  compiler to their CP-SAT model, guaranteeing constraint
  parity across pipelines.
\item \textbf{Zero-drift}.
  \texttt{seed\_implicit\_hardcoded(profs)} translates the
  legacy hardcoded HARDs into DSL. Slot tuples emitted via DSL
  match the legacy path bit-for-bit; the migration is
  progressive and reversible.
\item \textbf{HARD up-front, SOFT post-hoc}. DSL HARDs become
  CP-SAT constraints before the solve. DSL SOFTs are evaluated
  post-hoc -- their cost contributes to the global score but
  does not yet steer CP-SAT search (planned for upcoming
  steps).
\end{itemize}

\section{Backend lifecycle}

\texttt{webui/backend/main.py} defines a \texttt{lifespan}
context that calls \texttt{init\_db()} at startup. It does:

\begin{enumerate}
  \item \texttt{Base.metadata.create\_all(bind=engine)} -- create
        missing tables.
  \item \texttt{\_apply\_lightweight\_migrations()} -- idempotent
        ALTER TABLEs for fields added in later versions
        (\texttt{school\_classes.curriculum\_id},
        \texttt{teachers.last\_name}, \dots). Each ALTER is
        guarded by \texttt{PRAGMA table\_info} to avoid
        duplication.
\end{enumerate}

\medskip
\noindent
\textit{This chapter is the English edition. The Italian
edition (\texttt{docs/manual/chapters/architettura.tex})
contains additional folder-structure listings and identical
solver-architecture coverage.}
